\solution{Стабильность звезды}{10.0}

\tasksol Геометрически, угол $\Delta\varphi$ заметается за полгода, когда Земля смещена на $2R_E$. Тригонометрией получаем

\begin{eq}[points=0.2]
d = R_E \cot\left(\frac{\Delta\varphi}2\right)\approx\frac{2R_E}{\Delta\varphi}.
\end{eq}

Численно,

\begin{eq}[points=0.2]
d=\frac{2\times\num{1.5e11}}{\num{66.7e-3}\times\frac1{3600}\times\frac\pi{180}}=\qty{9.28e17}{\m}=\qty{98.0}{\ly}
\end{eq}

\criterion[type=penalty, points=-0.1]{Не записано в единицах \unit{\ly}}

\tasksol Всего мы имеем

\begin{eq}[points=0.2]
N = n\cdot \frac43\pi R^3
\end{eq}

молекул. Энергия каждой молекулы равна $\frac32 kT$ (мы учитываем только поступательные степени свободы для околонулевых температур). Поэтому мы имеем

\begin{eq}[points=0.2]
U = \frac32 NkT.
\end{eq}

Комбинируя эти два уравнения, получаем результат.

\tasksol В данной задаче достаточно вспомнить, что напряжённость гравитационного поля внутри однородного шара растёт линейно с радиусом (это можно показать либо через теорему Гаусса для гравитационного поля, либо через теорему Ньютона о сферической оболочке). Имея на поверхности

\begin{eq}[points=0.2]
g_0=\frac{GM}{R^2},
\end{eq}

мы получаем функцию

\begin{eq}[points=0.3]
g(r)= g_0\cdot\frac rR.
\end{eq}

\tasksol Взглянуть на эту задачу можно двумя различными способами. Более быстрый заключается в том, чтобы понять, что $|E_G|$ имеет смысл\footnote{В более общем случае говоря, $\max\{0, -E_G\}$, если бы мы работали с электростатической задачей.} минимальной работы, которую требуется затратить против сил притяжения, чтобы ``растаскать'' облако на бесконечность, чтобы энергия притяжения стала равной нулю. Допустим, облако было уменьшено до переменного радиуса $r$. Тогда малый сферический элемент $dm$ имеет потенциальную энергию

\begin{eq}[points=0.3]
\phi(r) = -\frac{GM(r)}r=-\frac{GM r^2}{R^3}.
\end{eq}

Тогда $E_G$  равен интегралу

\begin{eq}[points=0.2]
E_G = \int_0^R \phi(r)\,dm.
\end{eq}

Альтернативно, но чуть более детально, можно рассмотреть распределение гравитационного потенциала $\varphi(r)$ (т.е. потенциальная энергия на единицу массы) внутри шара, а затем просуммировать по всем элементам облака. С таким методом мы получаем\footnote{Тут следует быть аккуратнее. Более строго говоря, мы определяем $\varphi = -\frac{\partial g_r}{\partial r}$, однако $g_r=-g(r)$ с учётом того, что радиальная компонента в сферической системе координат направлена ``внаружу'' облака, а $g(r)$ --- ``внутрь''.}
$$
\varphi(r) - \varphi(R) = \int_R^r g(r)\, dr.
$$
Здесь
$$
\varphi(R) = -\frac{GM}R,
$$
так что применяя результат из \solref{3}, получаем
$$
\varphi(r) = -\frac{GM}R\ab(\frac32-\frac{r^2}{R^2}).
$$
Потенциальная энергия находится интегрированием
$$
E_G = \frac12\int_0^R \varphi(r)\,dm,
$$
где $dm=\rho\cdot 4\pi r^2\,dr= 3M r^2\,dr/R^3$, а коэффициент $1/2$ учитывает дважды просуммированные взаимодействия. Интегрируя, получаем результат

\begin{eq}[points=0.5]
E_G=-\frac35\frac{GM}R\iff f=\frac35.
\end{eq}

\tasksol Нужно добавить соотношение

\begin{eq}[points=0.2]
M = mN,
\end{eq}

где $N$ выражен в \eqref{3}. Полная энергия равна $E=E_G+U$, так что в критическом случае $|E_G|=U$. Сочетая результаты \eqref{3}--\eqref{4} и \eqref{9}, получаем

\begin{eq}[points=0.3]
    M_J=\frac{5k T}{4Gm^2}\sqrt\frac{15 k T}{2\pi nG}=\qty{2.60e31}{\kg}.
\end{eq}

\criterion[type=remark]{Записывать $M_J$ в кг необязательно}

В единицах массы Солнца,
\begin{eq}[points=0.1]
M_J=13 M_\odot.
\end{eq}

Понятное дело, это \textit{минимальная} масса. Коллапс произойдёт при любом $M>M_J$ при фиксированных прочих параметрах.

\tasksol Рассмотрим диск массой $dM=\rho(r)\cdot S dr$. Сверху на него действует сила давления $P(r+dr)S$, а снизу --- $P(r)S$. Вдобавок мы имеем направленный вниз $dM\cdot g(r)$, ввиду которого и создаётся градиент давлений. Уравнение гидростатического равновесия говорит

\begin{eq}[points=0.2]
P(r)S - P(r+dr)S-dM\cdot g(r)=0.
\end{eq}

Осталось выразить $g(r)$. Обобщая \eqref{6}, в общем виде мы запишем

\begin{eq}[points=0.2]
g(r)=\frac{GM(r)}{r^2}.
\end{eq}

Поскольку $P(r+dr)-P(r)=dP$, получаем

\begin{eq}[points=0.3]
\frac{dP}{dr}=-\frac{GM(r)\rho(r)}{r^2}.
\end{eq}


\tasksol Технически, нужно заново вывести формулу для $E_G$ в более общем виде, \textit{не} полагая $\rho(r)=\operatorname{const}$, что и будет отличаться от \eqref{9}. Используем \eqref{7} не раскрывая $M(r)$. Тогда мы получим тот же самый интеграл \eqref{8}, где распишем $dm=\rho(r)\cdot4\pi r^2\,dr$. Тогда можно заметить, что $r\cdot dP/dr=\phi(r)\rho(r)=\phi(r)\,dm/dV$, или же

\begin{eq}[points=0.4]
E_G = \int_0^{V_\odot} r\frac{dP}{dr}\,dV=\int_0^{R_\odot}4\pi r^3\,dP=\int_0^{R_\odot}-4\pi G M(r)\rho(r) r\,dr.
\end{eq}

\criterion[type=remark]{Принимается любая аналогичная запись, если $\rho(r)$ и $P(r)$ сохранены в общем виде}

В зависимости от предпочтения участника можно по-разному записать этот интеграл: три возможных примера записи даны выше, и балл даётся за любую верную запись интеграла, где \textit{не} было применено предположения $\rho(r)=\operatorname{const}$. Далее это записывается как

\begin{eq}[points=0.3]
E_G=\int 3V\,dP=3PV\Big|_0^{R_\odot}-\int_0^{V_\odot}3P\,dV=-3\int_0^{V_\odot}P\,dV.
\end{eq}

\criterion[type=remark]{Показанное интегрирование по частям принимается \textbf{только} если показано, что $PV\big|_0^{r_\odot}=0$}

где в первом `=' было применено интегрирование по частям, причём $PV|_0^{R_\odot}$ пропадает в нижнем пределе из-за равенства нулю объёма и в верхнем пределе из-за равенства нулю давления. В результате получаем

\begin{eq}[points=0.3]
E_G = -3\langle P\rangle V_\odot\iff \beta=-3.
\end{eq}

\tasksol Подставив $f=1$ в \eqref{9} и приравняв к \eqref{18}, мы сразу же получаем

\begin{eq}[points=0.3, numpoints=0.2]
\langle P\rangle = \frac{GM_\odot^2}{4\pi r^4_\odot}=\qty{9e13}{\Pa}.
\end{eq}

\tasksol Уравнение состояния идеального газа записывается в виде

\begin{eq}[points=0.2]
    \langle P\rangle = \frac{\langle\rho\rangle}{\overline m}k T_0 \iff \langle P\rangle V_\odot = \langle\nu\rangle R T_0,
\end{eq}

где $\langle\rho\rangle=\frac{M_\odot}{V_\odot}\iff\langle\nu\rangle=\frac{M_\odot}{\overline m N_A}$. В итоге

\begin{eq}[points=0.2, numpoints=0.1]
T_0=\frac{\overline m \langle P\rangle V_\odot}{k M_\odot}= \frac{GM_\odot\overline m}{3kR_\odot}=\qty{4.7e6}{\K}.
\end{eq}

\tasksol Стандартное применение закона Стефана--Больцмана. Записав

\begin{eq}[points=0.2]
    L_\odot=\sigma T^4_s\cdot4\pi R_\odot^2,
\end{eq}

находим

\begin{eq}[points=0.2, prefix={Численное значение}]
T_s=\ab(\frac{L_\odot}{4\pi\sigma R_\odot^2})^{1/4}\approx\qty{5800}{\K}.
\end{eq}

\tasksol Суммарное перемещение фотона равно

\begin{eq}[points=0.2]
    \vec D = \vec l_1+\vec l_2+\dots+\vec l_N.
\end{eq}

Поскольку перемещение фотона аппроксимировано как чисто стохастическое, все направления перемещения равновероятны, и потому в среднем

\begin{eq}[points=0.3]
    \langle\vec D\rangle=0.
\end{eq}

\tasksol Возьмём \eqref{24} в квадрат. Выйдет

\begin{eq}[points=0.3, override={D^2 = l_1^2+\dots+l_N^2+2\ab(\vec l_1\cdot\vec l_2+\vec l_1\cdot\vec l_3+\dots)}]
    D^2 = l_1^2+\dots+l_N^2+2\ab(\vec l_1\cdot\vec l_2+\vec l_1\cdot\vec l_3+\dots)=\sum_{i=1}^N l_i^2 + 2\sum_{i\neq j}\vec l_i\cdot\vec l_j,
\end{eq}

\criterion[type=penalty, points=-0.2]{Не показаны попарные суммы без объяснения}

где выражение в скобках принимает все попарные суммы по неодинаковым индексам. Учитывая, что $|l_i|=l\ \forall i$, то первая сумма даёт $Nl^2$. Снова учитывая стохастичность движения фотона, видим, что в среднем вторая сумма даёт нуль. В итоге

\begin{eq}[points=0.2]
    \langle D^2\rangle = Nl^2.
\end{eq}

Среднеквадратичное смещение покинувшего Солнце фотона должно быть приблизительно равно $r^2_\odot$. Итого фотону нужно порядка

\begin{eq}[points=0.2]
    \langle D^2\rangle\approx R_\odot^2\iff N=\ab(\frac{R_\odot}l)^2
\end{eq}

шагов до покидания Солнца, а это соответствует времени

\begin{eq}[points=0.1]
    t = N\cdot\frac lc.
\end{eq}

Итак,

\begin{eq}[points=0.2]
    t=\frac{R_\odot^2}{cl}.
\end{eq}

\tasksol Как было описано в условии, мы можем записать

\begin{eq}[points=0.3]
    L_\odot' = \sigma T_0^4\cdot 4\pi R_\odot^2.
\end{eq}

Связывая с \eqref{22} и информацией из условия, получаем

\begin{eq}[points=0.2, numpoints=0.2]
l = R_\odot\ab(\frac{T_s}{T_0})^4\approx\qty{1.67}{\mm}.
\end{eq}

Это показывает, что из \eqref{30} время численно равно

\begin{eq}[points=0.3, prefix={Численное значение}]
    t=\qty{9.67e11}{\s}\approx\num{30600}~\text{лет}.
\end{eq}

\tasksol Имея уравнение \eqref{21}, подставим его в выражение \eqref{31}. Домножив \eqref{31} на $l/R_\odot$, это будет выражением для $L_\odot$. Тогда выйдет

\begin{eq}[points=0.3]
    L_\odot = \frac{4\pi \sigma G^4 m^4 l}{81 k^4}\cdot\frac{M_\odot^4}{R_\odot^3}.
\end{eq}

Вытащив $\langle\rho\rangle$ в этом выражении, заключаем

\begin{eq}[points=0.2]
    C=\frac{16 \pi^2 \sigma G^4 m^4 l}{243 k^4},
\end{eq}

\begin{eq}[points=0.2]
n=3.
\end{eq}

\tasksol Из последнего уравнения мы имеем

\begin{equation}
\log L = \log L + n\cdot \log M,
\end{equation}

так что в координатах $\log L$--$\log M$ мы аппроксимируем заданные в графике по условию точки прямой линией. Следует обратить внимание, что ось икс хоть и имеет логарифмическую шкалу, нужно пересчитать значения на оси, взяв логарифм. Тогда получаем так, как по рисунку ниже.

\cpic{1}{figures/2.6.png}

Возьмём точки $(1.3,5.4)$ и $(-0.7,-2.5)$ у сглаживающей прямой. Тогда коэффициент наклона равен
$$
n=\frac{(-2.5) - (+5.4)}{(-0.7) - (+1.3)}=3.95.
$$
Результат немногим больше полученного в \solref{14} ответа.

\criterion[points=0.8]{Ответ в интервале $n\in[3.8,\ 4.1]$}
\criterion[type=penalty, points=-0.4]{Ответ в интервале $n\in[3.5,\ 4.4]$}